% BEMv19 appendix snippet
% Multi-component link geometry parser for Route-B.

\subsection{Multi-component link geometry parser}

\paragraph{Status.}
This appendix introduces the geometry parser required before links can be
included in the Route-B R--T spectral programme.  It does not yet define the
full multi-component R--T boundary operator.

For a link \(L=C_1\cup\cdots\cup C_m\), each component is represented by its
own Fourier curve
\[
\mathbf x_a(t)=
\sum_{n\ge1}
\left[
\mathbf A_{a,n}\cos(nt)+\mathbf B_{a,n}\sin(nt)
\right],
\qquad a=1,\ldots,m .
\]
The parser evaluates each component, computes its arclength, and estimates
pairwise Gauss linking numbers
\[
\operatorname{Lk}(C_a,C_b)
=
\frac{1}{4\pi}
\oint_{C_a}\oint_{C_b}
\frac{(\mathbf x-\mathbf y)\cdot(d\mathbf x\times d\mathbf y)}
{\|\mathbf x-\mathbf y\|^3}.
\]

\paragraph{Route-B operator upgrade.}
The single-component boundary map must be replaced by a block operator
\[
\Lambda^{\rm link}_{R/T}
=
\left(\Lambda^{R/T}_{ab}\right)_{a,b=1}^{m},
\]
where diagonal blocks describe self-response and off-diagonal blocks describe
inter-component R--T coupling.  Only after this block operator is defined can
links enter the same spectral determinant used for knots.

\paragraph{Status tag.}
\[
\boxed{\text{BEMv19: link geometry parser ready; block R--T operator open.}}
\]
